\chapter{La señal de alarma}
% versión completa: chapters/14_pain.tex

\pencilfig{14}{\pencilart{14}}{El dolor es una alarma de incendio, no un termómetro.}

Todos los demás sentidos informan de cómo es el mundo. El dolor, no. Informa de que algo amenaza al cuerpo, y su volumen lo fija en gran parte la propia máquina, no el estímulo. El dolor es una alarma de incendio, no un termómetro.

Los sensores del dolor solo se activan ante lo fuerte: calor de más de cuarenta y tantos grados, presión fuerte, sustancias corrosivas. Se acostumbran mucho menos que los sensores del tacto: la alarma no debe callarse mientras haya amenaza.

\section*{Dos cables}

Recuerde cuando se golpeó un dedo. Primero, un dolor agudo y preciso: dónde y cuándo. Un instante después, uno sordo, difuso, largo: qué tan grave. Son dos cables distintos. Uno es rápido: su señal llega en decenas de milisegundos. El otro es lento: su señal tarda del orden de un segundo.

\section*{La compuerta}

¿Por qué, tras un golpe, nos frotamos la zona? En la médula espinal, la señal de dolor pasa por una “compuerta”. El tacto despierta una neurona inhibidora que la entorna. Por eso frotar el golpe de verdad ayuda, pero solo contra un dolor moderado; el fuerte pasa igual. Que el dolor fuerte apague a su vez la neurona inhibidora es parte de la teoría original de la compuerta, y todavía no tiene una confirmación directa.

\pencilfig{126}{%
% gate: the pain wire drives the relay neuron; touch wakes an inhibitory neuron that closes the gate
\draw[pen=pcAmber] (0,1.35) -- (3.05,1.0); \draw[arr=pcAmber] (2.8,1.03) -- (3.08,1.0);
\node[lbl, text=pcAmber!70!black, anchor=east] at (-0.05,1.35) {dolor};
\draw[pen=pcBlue] (0,-0.15) -- (1.75,-0.15); \draw[arr=pcBlue] (1.5,-0.15) -- (1.78,-0.15);
\node[lbl, text=pcBlue, anchor=east] at (-0.05,-0.15) {tacto};
\draw[dotpen=pcRed, wash=pcRed] (2.05,-0.15) circle (0.26);
\draw[pcRed, line width=0.7pt, -{Bar[width=6pt]}] (2.25,0.05) -- (3.2,0.66);
\node[lbl, text=pcRed, anchor=north] at (2.05,-0.45) {freno};
\draw[dotpen=pcGray, wash=pcGray] (3.4,0.9) circle (0.34);
\node[lbl, anchor=south] at (3.4,1.3) {compuerta};
\draw[pen] (3.75,0.9) -- (5.6,0.9); \draw[arr] (5.35,0.9) -- (5.65,0.9);
\node[lbl, anchor=west] at (5.7,0.9) {al cerebro};
}{La compuerta de la médula espinal. La señal de dolor pasa al cerebro a través de una neurona compuerta, y el tacto despierta una neurona inhibidora que la entorna. Por eso frotar el golpe ayuda, pero solo contra un dolor moderado.}

\section*{Una perilla de volumen desde arriba}

El cerebro sabe subir o bajar el volumen del dolor por su cuenta. Desde arriba, desde el encéfalo, bajan hasta la compuerta canales de radio: \nt{SERO}, \nt{NORA} y unas sustancias especiales, parecidas a los analgésicos, que el cerebro produce él mismo. En una situación de peligro el dolor puede casi desaparecer: quien huye del peligro no tiene tiempo de cojear. La expectativa de alivio también atenúa el dolor, en parte a través de esas mismas sustancias.

\section*{Cuando la alarma se traba}

El dolor repetido aumenta la sensibilidad: el mismo golpe empieza a doler más. También es aprendizaje: un lazo que se refuerza a sí mismo mientras la herida sana. Pero un lazo fuerte tiene un peligro: puede “trabarse” y seguir sonando cuando la herida ya sanó, como el grupo de neuronas del segundo capítulo que se mantiene solo.

\section*{Maestro e interrupción}

El dolor es además un maestro: para el “mejor o peor de lo esperado” de la dopamina funciona como una recompensa negativa. Y también interrumpe la tarea en curso, como un estallido de noradrenalina: deja todo y ocúpate del peligro. Y la reacción más rápida, retirar la mano de algo caliente, se cierra en la propia médula espinal con ese mismo par de músculos opuestos del capítulo anterior, antes incluso de que usted llegue a sentir el dolor.

\fullversion{14}
